% One wire for the event input, and the instrument this bench does not have.
\begin{tikzpicture}[font=\sffamily\small, x=1cm, y=1cm]

\node[board, minimum width=46mm, minimum height=26mm] (nucleo) at (0,0) {};
\node[text=white, font=\sffamily\bfseries\small] at (0,0.62) {NUCLEO-H7A3ZI-Q};
\node[text=white, font=\sffamily\scriptsize] at (0,0.18) {TIM2 CH1, input capture};
\node[pinrow, minimum width=40mm] at (0,1.05) {};
\node[ledsym, fill=green!65!black] at (-1.5,-0.55) {};
\node[text=white, font=\sffamily\tiny, anchor=west] at (-1.25,-0.55) {32.768 kHz crystal, fitted};

\node[ext, text width=30mm] (src) at (5.4,0.55)
  {an edge to stamp\\{\scriptsize the user button, or a line from the profiler}};
\node[absentblk, text width=30mm] (gps) at (5.4,-1.6)
  {receiver with a pulse per second\\{\scriptsize not on this bench}};

\draw[cable, draw=red!60!black] (nucleo.east) -- node[lbl, above]{one wire} (src.west);
\draw[hiz, ->] (gps.west) -- ++(-1.2,0) |- (nucleo.east);

\node[ext, text width=26mm] (host) at (-5.2,0) {host PC\\{\scriptsize console}};
\draw[usbcable] (host.east) -- node[lbl, above]{USB} (nucleo.west);

\node[umlnote, text width=58mm, anchor=north] at (0.4,-2.4)
  {What the wire gives: an externally generated edge whose timestamp is taken by
   the timer at the pin. The difference between that stamp and one taken at the
   top of the interrupt handler is this board's interrupt latency, measured
   rather than quoted.\\[3pt]
   What the absent part would give: an absolute reference. The drift figure this
   chapter reports is a comparison of the board's two oscillators and cannot say
   which of them is wrong. A receiver emitting one pulse per second, captured on
   this same channel, would turn that comparison into a frequency measurement
   with a stated uncertainty. It is the cheapest instrument that would improve
   this volume, and it is named here rather than quietly wished for.};
\end{tikzpicture}
